\chapter{Nhóm bài toán ghép nhóm và phối hợp điều kiện}

\section{Bài toán Qua sông (River Crossing Problem / Hackers and Serfs)}

Bài toán yêu cầu ghép một nhóm người lên thuyền theo các tổ hợp hợp lệ. Ví dụ biến thể Hackers and Serfs thường yêu cầu mỗi thuyền có bốn người, nhưng không được tạo ra tổ hợp gây mất cân bằng không hợp lệ.

\subsection{Mô tả bài toán}
Có hai loại tiến trình: \textit{hacker} và \textit{serf}. Mỗi thuyền chở đúng 4 người; chỉ các tổ hợp sau được phép khởi hành:
\begin{itemize}
    \item 4 hacker (4H),
    \item 4 serf (4S),
    \item 2 hacker + 2 serf (2H2S).
\end{itemize}
Mục tiêu là ghép nhóm đúng điều kiện, không để một nhóm không hợp lệ lên thuyền, và đảm bảo chỉ khi đủ 4 người thì thuyền mới rời bến. Bài toán được trình bày phổ biến trong \cite{downey2016}.

\subsection{Các thực thể tham gia}
\begin{itemize}
    \item \textbf{Hacker:} tiến trình đại diện cho người thuộc nhóm hacker, đến bến và chờ được ghép vào một thuyền hợp lệ.
    \item \textbf{Serf:} tiến trình đại diện cho người thuộc nhóm serf, có hành vi tương tự nhưng thuộc loại khác để kiểm tra ràng buộc tổ hợp.
    \item \textbf{Captain:} một tiến trình trong nhóm bốn người được chọn để thực hiện thao tác khởi hành sau khi barrier xác nhận cả nhóm đã lên thuyền.
\end{itemize}

\subsection{Tài nguyên dùng chung}
\begin{itemize}
    \item Các biến đếm số người đang chờ: \texttt{waitingH}, \texttt{waitingS}.
    \item Trạng thái ghép nhóm/đánh thức (quyết định ai được lên thuyền ở lượt hiện tại).
\end{itemize}
Các biến đếm phải được bảo vệ bởi mutex để tránh race condition khi nhiều người đến cùng lúc.

\subsection{Ràng buộc đồng bộ}
\begin{itemize}
    \item Một tiến trình chỉ được \textit{board} (lên thuyền) khi nó thuộc một nhóm hợp lệ vừa được chọn.
    \item Đúng 4 tiến trình phải đồng bộ với nhau: không ai được “khởi hành” khi chưa đủ 4 người.
    \item Không tiến trình nào bị đánh thức nhầm (không để 5 người cùng lên cho một lượt).
\end{itemize}

\subsection{Phân tích lỗi tiềm ẩn}

\subsubsection*{Race condition}
Nếu quên bảo vệ các biến đếm bằng mutex, nhiều tiến trình đến cùng lúc có thể
cùng nhìn thấy một trạng thái cũ và chọn sai tổ hợp, dẫn đến việc cho phép một
nhóm không hợp lệ lên thuyền.

\subsubsection*{Deadlock}
Nguy cơ deadlock xuất hiện khi tiến trình giữ \texttt{mutex} trong lúc chờ trên
\texttt{hQueue}, \texttt{sQueue} hoặc barrier. Lời giải nhả mutex trước khi các
tiến trình không phải captain đi ngủ, còn captain chỉ giữ quyền điều phối cho
đến khi nhóm hiện tại hoàn tất.

\subsubsection*{Starvation / fairness}
Nếu dòng người đến quá lệch, một loại tiến trình có thể chờ lâu. Có thể bổ
sung chính sách ưu tiên tổ hợp 2H2S khi khả dụng để giảm lệch, nhưng lời giải
cốt lõi chỉ bảo đảm đúng tổ hợp và không cho người thứ năm chen vào cùng
chuyến.

\subsection{Giải pháp đề xuất}
Ý tưởng chuẩn là dùng:
\begin{itemize}
    \item 1 mutex để bảo vệ \texttt{waitingH}, \texttt{waitingS} và thao tác chọn nhóm.
    \item 2 semaphore hàng đợi: \texttt{hQueue} và \texttt{sQueue} để đánh thức đúng người được chọn.
    \item 1 cơ chế \textit{barrier} cho nhóm 4 người (có thể cài bằng semaphore + bộ đếm) để đảm bảo cả 4 người đã lên thuyền trước khi “khởi hành”.
\end{itemize}
Trong mỗi nhóm 4 người, chọn 1 tiến trình đóng vai trò \textit{captain} thực hiện hành động “row/khởi hành”.

\subsection{Mã giả}

\captionof{algorithm}{River Crossing (Hackers and Serfs) -- giả mã dùng semaphore}
\label{alg:river-crossing}
\begin{algorithmic}[1]
\State \textbf{Shared:} \texttt{mutex} $\leftarrow 1$; \texttt{waitingH} $\leftarrow 0$; \texttt{waitingS} $\leftarrow 0$
\State \textbf{Shared:} \texttt{hQueue} $\leftarrow 0$; \texttt{sQueue} $\leftarrow 0$
\State \textbf{Shared (barrier 4):} \texttt{bMutex} $\leftarrow 1$; \texttt{bCount} $\leftarrow 0$; \texttt{bSem1} $\leftarrow 0$; \texttt{bSem2} $\leftarrow 0$

\Function{Barrier4}{}
    \Comment{Pha 1: đợi đủ 4 người đã được chọn lên thuyền}
    \State \Call{wait}{\texttt{bMutex}}
    \State $\texttt{bCount} \gets \texttt{bCount} + 1$
    \If{$\texttt{bCount} = 4$}
        \For{$i \gets 1$ \textbf{to} $4$}
            \State \Call{signal}{\texttt{bSem1}}
        \EndFor
    \EndIf
    \State \Call{signal}{\texttt{bMutex}}
    \State \Call{wait}{\texttt{bSem1}}

    \Comment{Pha 2: đợi cả 4 người rời barrier trước khi tái sử dụng}
    \State \Call{wait}{\texttt{bMutex}}
    \State $\texttt{bCount} \gets \texttt{bCount} - 1$
    \If{$\texttt{bCount} = 0$}
        \For{$i \gets 1$ \textbf{to} $4$}
            \State \Call{signal}{\texttt{bSem2}}
        \EndFor
    \EndIf
    \State \Call{signal}{\texttt{bMutex}}
    \State \Call{wait}{\texttt{bSem2}}
\EndFunction

\Procedure{Hacker}{}
    \State $\texttt{captain} \gets false$
    \State \Call{wait}{\texttt{mutex}}
    \State $\texttt{waitingH} \gets \texttt{waitingH} + 1$
    \If{$\texttt{waitingH} = 4$}
        \For{$i \gets 1$ \textbf{to} $4$}
            \State \Call{signal}{\texttt{hQueue}}
        \EndFor
        \State $\texttt{waitingH} \gets \texttt{waitingH} - 4$
        \State $\texttt{captain} \gets true$
    \ElsIf{$\texttt{waitingH} = 2$ \textbf{and} $\texttt{waitingS} \ge 2$}
        \For{$i \gets 1$ \textbf{to} $2$}
            \State \Call{signal}{\texttt{hQueue}}
            \State \Call{signal}{\texttt{sQueue}}
        \EndFor
        \State $\texttt{waitingH} \gets \texttt{waitingH} - 2$; $\texttt{waitingS} \gets \texttt{waitingS} - 2$
        \State $\texttt{captain} \gets true$
    \Else
        \State \Call{signal}{\texttt{mutex}} \Comment{không phải captain thì nhả mutex trước khi chờ}
    \EndIf

    \State \Call{wait}{\texttt{hQueue}} \Comment{chỉ được lên thuyền khi được chọn}
    \State \Call{Board}{}
    \State \Call{Barrier4}{}
    \If{$\texttt{captain}$}
        \State \Call{RowBoat}{}
        \State \Call{signal}{\texttt{mutex}} \Comment{mo khoa sau khi chuyen da xong}
    \EndIf
\EndProcedure

\Procedure{Serf}{}
    \State $\texttt{captain} \gets false$
    \State \Call{wait}{\texttt{mutex}}
    \State $\texttt{waitingS} \gets \texttt{waitingS} + 1$
    \If{$\texttt{waitingS} = 4$}
        \For{$i \gets 1$ \textbf{to} $4$}
            \State \Call{signal}{\texttt{sQueue}}
        \EndFor
        \State $\texttt{waitingS} \gets \texttt{waitingS} - 4$
        \State $\texttt{captain} \gets true$
    \ElsIf{$\texttt{waitingS} = 2$ \textbf{and} $\texttt{waitingH} \ge 2$}
        \For{$i \gets 1$ \textbf{to} $2$}
            \State \Call{signal}{\texttt{sQueue}}
            \State \Call{signal}{\texttt{hQueue}}
        \EndFor
        \State $\texttt{waitingS} \gets \texttt{waitingS} - 2$; $\texttt{waitingH} \gets \texttt{waitingH} - 2$
        \State $\texttt{captain} \gets true$
    \Else
        \State \Call{signal}{\texttt{mutex}} \Comment{khong phai captain thi nho mutex truoc khi cho}
    \EndIf

    \State \Call{wait}{\texttt{sQueue}}
    \State \Call{Board}{}
    \State \Call{Barrier4}{}
    \If{$\texttt{captain}$}
        \State \Call{RowBoat}{}
        \State \Call{signal}{\texttt{mutex}} \Comment{mở khóa sau khi chuyến đã xong}
    \EndIf
\EndProcedure
\end{algorithmic}

\subsection{Kiểm chứng ngắn}

Nếu bến có đủ 2 hacker và 2 serf, tiến trình cuối cùng làm điều kiện đủ sẽ phát đúng 2 tín hiệu cho \texttt{hQueue} và 2 tín hiệu cho \texttt{sQueue}. Bốn tiến trình này cùng đi qua \texttt{Barrier4}; chỉ khi cả bốn người đã lên thuyền thì captain mới gọi \texttt{RowBoat}. Trong mọi lượt, số tín hiệu phát ra đúng bằng số vị trí trên thuyền, nên không thể có người thứ năm chen vào cùng chuyến.

\subsection{Kết luận bài toán}

River Crossing minh họa đồng bộ theo điều kiện ghép nhóm. Mutex bảo vệ biến đếm, hai semaphore hàng đợi đánh thức đúng loại tiến trình, còn barrier bảo đảm nhóm đủ bốn người trước khi khởi hành. Lời giải đúng cần giữ bất biến mỗi chuyến chỉ thuộc một trong ba tổ hợp hợp lệ: 4H, 4S hoặc 2H2S.

\section{Bài toán Những người hút thuốc (Cigarette Smokers Problem)}

Bài toán gồm một tác nhân đặt hai trong ba nguyên liệu lên bàn, và người hút thuốc có nguyên liệu còn lại sẽ được hút. Mục tiêu là đánh thức đúng tiến trình phù hợp với cặp nguyên liệu hiện có.

\subsection{Mô tả bài toán}
Có 3 loại nguyên liệu: thuốc lá (T), giấy (P) và diêm (M). Có 1 \textit{agent} mỗi lượt đặt lên bàn đúng 2 nguyên liệu. Có 3 \textit{smoker}: mỗi smoker có sẵn một nguyên liệu và chỉ cần 2 nguyên liệu còn lại để cuốn thuốc và hút.

Mục tiêu: ở mỗi lượt, \textbf{chỉ} smoker phù hợp với cặp nguyên liệu trên bàn được phép tiếp tục; các smoker còn lại phải chờ. Bài toán và lời giải kinh điển bằng semaphore được trình bày trong \cite{downey2016}.

\subsection{Các thực thể tham gia}
\begin{itemize}
    \item \textbf{Agent:} tiến trình đặt đúng hai nguyên liệu lên bàn ở mỗi lượt.
    \item \textbf{Pusher:} các tiến trình trung gian ghi nhận nguyên liệu vừa xuất hiện và ghép cặp an toàn dưới mutex.
    \item \textbf{Smoker:} ba tiến trình có sẵn lần lượt thuốc lá, giấy hoặc diêm; mỗi tiến trình chỉ được chạy khi hai nguyên liệu còn lại đã sẵn sàng.
\end{itemize}

\subsection{Tài nguyên dùng chung}
\begin{itemize}
    \item Trạng thái “bàn” (có những nguyên liệu nào đang có sẵn).
    \item Quyền được lấy nguyên liệu/hút ở mỗi lượt (tránh nhiều smoker cùng lấy).
\end{itemize}

\subsection{Ràng buộc đồng bộ}
\begin{itemize}
    \item Agent chỉ được đặt nguyên liệu khi lượt trước đã kết thúc (một smoker đã hút xong).
    \item Khi agent đặt (T,P) thì chỉ smoker có M được đánh thức; tương tự các cặp còn lại.
    \item Không được bỏ lỡ tín hiệu (missed wakeup) khi hai nguyên liệu đến “lệch nhịp”.
\end{itemize}

\subsection{Phân tích lỗi tiềm ẩn}

\subsubsection*{Race condition}
Nếu các pusher cập nhật cờ \texttt{hasT}, \texttt{hasP}, \texttt{hasM} mà
không khóa, hai pusher có thể cùng quyết định trên trạng thái bàn không nhất
quán và đánh thức sai smoker.

\subsubsection*{Deadlock}
Agent phải chờ \texttt{agentSem} trước khi đặt lượt mới và smoker phải
\texttt{signal(agentSem)} sau khi hút xong. Nếu thiếu tín hiệu này, agent có
thể ngủ vô hạn dù lượt trước đã kết thúc.

\subsubsection*{Starvation / fairness}
Nếu agent chọn cặp nguyên liệu lệch trong thời gian dài, smoker tương ứng có
thể chờ lâu. Đây là đặc tính của nguồn vào; phần đồng bộ chỉ bảo đảm mỗi lượt
có đúng một smoker phù hợp được đánh thức.

\subsection{Giải pháp đề xuất}
Lời giải tiêu chuẩn dùng mô hình \textit{pushers}:
\begin{itemize}
    \item Agent phát tín hiệu cho từng nguyên liệu (3 semaphore T/P/M).
    \item Ba tiến trình pusher quan sát nguyên liệu, ghép cặp dưới mutex thông qua các cờ $hasT$, $hasP$, $hasM$.
    \item Khi đủ 2 nguyên liệu, pusher phát tín hiệu cho đúng smoker (3 semaphore cho 3 smoker).
    \item Smoker hút xong sẽ \textit{signal} cho agent để bắt đầu lượt kế tiếp.
\end{itemize}

\subsection{Mã giả}

\captionof{algorithm}{Cigarette Smokers -- lời giải pushers (giả mã)}
\label{alg:cigarette-smokers}
\begin{algorithmic}[1]
\State \textbf{Shared:} \texttt{agentSem} $\leftarrow 1$
\State \textbf{Shared:} \texttt{tobacco} $\leftarrow 0$; \texttt{paper} $\leftarrow 0$; \texttt{match} $\leftarrow 0$
\State \textbf{Shared:} \texttt{smokerT} $\leftarrow 0$; \texttt{smokerP} $\leftarrow 0$; \texttt{smokerM} $\leftarrow 0$
\State \textbf{Shared:} \texttt{mutex} $\leftarrow 1$; \texttt{hasT} $\leftarrow false$; \texttt{hasP} $\leftarrow false$; \texttt{hasM} $\leftarrow false$

\Procedure{Agent}{}
    \While{true}
        \State \Call{wait}{\texttt{agentSem}}
        \State \Call{PlaceTwoIngredients}{} \Comment{ngẫu nhiên 2 trong 3}
        \If{placed (T,P)} \State \Call{signal}{\texttt{tobacco}}; \Call{signal}{\texttt{paper}} \EndIf
        \If{placed (T,M)} \State \Call{signal}{\texttt{tobacco}}; \Call{signal}{\texttt{match}} \EndIf
        \If{placed (P,M)} \State \Call{signal}{\texttt{paper}}; \Call{signal}{\texttt{match}} \EndIf
    \EndWhile
\EndProcedure

\Procedure{PusherT}{}
    \While{true}
        \State \Call{wait}{\texttt{tobacco}}
        \State \Call{wait}{\texttt{mutex}}
        \If{$\texttt{hasP}$}
            \State $\texttt{hasP} \gets false$; \Call{signal}{\texttt{smokerM}}
        \ElsIf{$\texttt{hasM}$}
            \State $\texttt{hasM} \gets false$; \Call{signal}{\texttt{smokerP}}
        \Else
            \State $\texttt{hasT} \gets true$
        \EndIf
        \State \Call{signal}{\texttt{mutex}}
    \EndWhile
\EndProcedure

\Procedure{PusherP}{}
    \While{true}
        \State \Call{wait}{\texttt{paper}}
        \State \Call{wait}{\texttt{mutex}}
        \If{$\texttt{hasT}$}
            \State $\texttt{hasT} \gets false$; \Call{signal}{\texttt{smokerM}}
        \ElsIf{$\texttt{hasM}$}
            \State $\texttt{hasM} \gets false$; \Call{signal}{\texttt{smokerT}}
        \Else
            \State $\texttt{hasP} \gets true$
        \EndIf
        \State \Call{signal}{\texttt{mutex}}
    \EndWhile
\EndProcedure

\Procedure{PusherM}{}
    \While{true}
        \State \Call{wait}{\texttt{match}}
        \State \Call{wait}{\texttt{mutex}}
        \If{$\texttt{hasT}$}
            \State $\texttt{hasT} \gets false$; \Call{signal}{\texttt{smokerP}}
        \ElsIf{$\texttt{hasP}$}
            \State $\texttt{hasP} \gets false$; \Call{signal}{\texttt{smokerT}}
        \Else
            \State $\texttt{hasM} \gets true$
        \EndIf
        \State \Call{signal}{\texttt{mutex}}
    \EndWhile
\EndProcedure

\Procedure{SmokerWithT}{}
    \While{true}
        \State \Call{wait}{\texttt{smokerT}} \Comment{có sẵn T, chờ P+M}
        \State \Call{MakeCigaretteAndSmoke}{}
        \State \Call{signal}{\texttt{agentSem}}
    \EndWhile
\EndProcedure

\Procedure{SmokerWithP}{}
    \While{true}
        \State \Call{wait}{\texttt{smokerP}} \Comment{có sẵn P, chờ T+M}
        \State \Call{MakeCigaretteAndSmoke}{}
        \State \Call{signal}{\texttt{agentSem}}
    \EndWhile
\EndProcedure

\Procedure{SmokerWithM}{}
    \While{true}
        \State \Call{wait}{\texttt{smokerM}} \Comment{có sẵn M, chờ T+P}
        \State \Call{MakeCigaretteAndSmoke}{}
        \State \Call{signal}{\texttt{agentSem}}
    \EndWhile
\EndProcedure
\end{algorithmic}

\subsection{Kiểm chứng ngắn}

Khi agent đặt hai nguyên liệu, mỗi nguyên liệu chỉ đánh thức pusher tương ứng. Các pusher cập nhật các cờ \texttt{hasT}, \texttt{hasP}, \texttt{hasM} trong vùng khóa, nên chỉ một pusher có thể quyết định cặp nguyên liệu hoàn chỉnh. Pusher đó phát đúng một tín hiệu đến smoker còn thiếu nguyên liệu thứ ba; sau khi smoker hút xong, \texttt{agentSem} mới được mở cho lượt tiếp theo.

\subsection{Kết luận bài toán}

Cigarette Smokers cho thấy việc đánh thức trực tiếp theo sự kiện đơn lẻ có thể gây lost wakeup nếu không lưu trạng thái trung gian. Mẫu pusher tách việc ghi nhận nguyên liệu khỏi việc đánh thức smoker, nhờ đó bảo đảm mỗi lượt chỉ có đúng một smoker phù hợp được chạy và agent không ghi đè lên lượt chưa hoàn tất.

\newpage
\section{Bài toán Ông già Noel (Santa Claus Problem)}
% =========================================================
% BAI TOAN ONG GIA NOEL -- SANTA CLAUS PROBLEM
% Nguon giai thuat:
% Allen B. Downey, The Little Book of Semaphores, 2nd ed.,
% Section 5.5: The Santa Claus Problem, pp. 139--142.
% =========================================================

\subsection{Mô tả bài toán}

\subsubsection*{Cốt truyện và bối cảnh}

Bài toán Ông già Noel (\textit{Santa Claus Problem}) là một bài toán
đồng bộ tiến trình kinh điển. Trong bài toán, Santa Claus đang ngủ tại
Bắc Cực và chỉ được đánh thức khi xảy ra một trong hai tình huống:

\begin{itemize}
\item Cả 9 con tuần lộc đã quay về sau kỳ nghỉ và sẵn sàng kéo xe
trượt tuyết đi phát quà.
\item Có một nhóm gồm đúng 3 chú lùn cùng gặp khó khăn trong quá
trình làm đồ chơi và cần Santa trợ giúp.
\end{itemize}

Nếu Santa thức dậy và nhận thấy đồng thời có đủ 9 tuần lộc đang chờ
cùng với một nhóm 3 chú lùn cần giúp đỡ, Santa phải ưu tiên chuẩn bị xe
trượt tuyết cho tuần lộc trước, vì việc phát quà quan trọng hơn \cite[Sec.~5.5, p.~137]{downey2016}.

Mục tiêu của hệ thống là bảo đảm Santa chỉ được đánh thức khi đủ điều
kiện, các tuần lộc được xử lý theo nhóm 9 con, các chú lùn được xử lý
theo nhóm 3 người, và không có nhóm chú lùn mới chen vào trước khi nhóm
hiện tại đã thực hiện việc nhận trợ giúp.

\subsection{Các thực thể tham gia}

\paragraph{Santa Claus.} Santa là tiến trình trung tâm. Santa ngủ khi chưa có yêu cầu hợp lệ.
Khi được đánh thức, Santa kiểm tra xem có đủ tuần lộc hay không trước;
nếu không, Santa kiểm tra nhóm chú lùn.

\begin{center}
\texttt{Ngủ $\rightarrow$ Được đánh thức $\rightarrow$ Kiểm tra điều kiện
$\rightarrow$ Xử lý công việc $\rightarrow$ Ngủ lại}
\end{center}

\paragraph{Tuần lộc (Reindeer).}
Có 9 tiến trình tuần lộc. Mỗi tuần lộc quay về Bắc Cực, cập nhật số
tuần lộc đang chờ và sau đó chờ Santa chuẩn bị xe. Tuần lộc thứ 9 là
tiến trình đánh thức Santa.

\begin{center}
\texttt{Đi nghỉ $\rightarrow$ Quay về $\rightarrow$ Chờ đủ nhóm 9 con
$\rightarrow$ Được mắc vào xe}
\end{center}

\paragraph{Chú lùn (Elves).}
Có nhiều tiến trình chú lùn. Khi một chú lùn gặp khó khăn, chú lùn đó
xin tham gia vào nhóm hiện tại. Chỉ một nhóm gồm đúng 3 chú lùn được
hình thành tại một thời điểm.

\begin{center}
\texttt{Làm đồ chơi $\rightarrow$ Gặp khó khăn $\rightarrow$ Gia nhập nhóm
3 người $\rightarrow$ Nhận trợ giúp $\rightarrow$ Làm việc tiếp}
\end{center}

\subsection{Tài nguyên dùng chung}

Các tài nguyên dùng chung trong bài toán gồm:

\begin{itemize}
\item Biến \texttt{reindeer}: số tuần lộc đã quay về và đang chờ.
\item Biến \texttt{elves}: số chú lùn đang thuộc nhóm xin trợ giúp
hiện tại.
\item Santa, vì tại một thời điểm Santa chỉ có thể xử lý một loại
yêu cầu.
\item Tín hiệu đánh thức Santa.
\item Quyền gia nhập nhóm chú lùn hiện tại.
\item Tín hiệu cho phép tuần lộc tiếp tục sau khi Santa chuẩn bị xe.
\end{itemize}

Hai biến \texttt{reindeer} và \texttt{elves} phải được bảo vệ, vì nhiều
tiến trình có thể cùng muốn cập nhật chúng.

\subsubsection*{Sơ đồ minh họa}

\begin{figure}[H]
\centering
\fbox{
\begin{minipage}{0.83\textwidth}
\centering
\vspace{0.45cm}
        \textbf{9 Tuần lộc}
        $\quad \longrightarrow \quad$
        \textbf{Santa Claus}
        $\quad \longrightarrow \quad$
        \textbf{Chuẩn bị xe trượt tuyết}

        \vspace{0.6cm}

        \textbf{3 Chú lùn}
        $\quad \longrightarrow \quad$
        \textbf{Santa Claus}
        $\quad \longrightarrow \quad$
        \textbf{Trợ giúp làm đồ chơi}

        \vspace{0.6cm}

        \textit{Nếu cả hai nhóm cùng sẵn sàng, Santa ưu tiên tuần lộc.}

        \vspace{0.45cm}
    \end{minipage}
}
\caption{Mô hình các thực thể trong bài toán Santa Claus}
\label{fig:santa-model}

\end{figure}

% =========================================================

\subsection{Ràng buộc đồng bộ}

\subsubsection*{Loại trừ tương hỗ}

Đoạn găng (\textit{critical section}) của bài toán nằm ở các thao tác
đọc và cập nhật hai biến đếm dùng chung:

\begin{lstlisting}
reindeer = reindeer + 1;
reindeer = reindeer - 9;

elves = elves + 1;
elves = elves - 1;
\end{lstlisting}

Các thao tác này không được thực hiện đồng thời mà không có khóa bảo vệ.
Ví dụ, nếu hai chú lùn cùng đọc và tăng biến \texttt{elves}, hệ thống
có thể ghi nhận sai số chú lùn đang chờ. Tương tự, nếu nhiều tuần lộc
cùng cập nhật \texttt{reindeer}, Santa có thể không được đánh thức đúng
thời điểm.

Vì vậy, giải pháp sử dụng semaphore \texttt{mutex} như một khóa loại
trừ tương hỗ để bảo vệ hai biến trạng thái.

\subsubsection*{Đồng bộ điều kiện}

\paragraph{Quan hệ chờ -- đánh thức của Santa.}

\begin{itemize}
\item Santa chờ trên semaphore \texttt{santaSem} khi chưa có yêu cầu
hợp lệ.
\item Tuần lộc thứ 9 gọi \texttt{signal(santaSem)} để đánh thức Santa.
\item Chú lùn thứ 3 trong nhóm gọi \texttt{signal(santaSem)} để đánh
thức Santa.
\end{itemize}

\paragraph{Quan hệ chờ -- đánh thức của tuần lộc.}

\begin{itemize}
\item Sau khi quay về, mỗi tuần lộc chờ trên
\texttt{reindeerSem}.
\item Khi có đủ 9 tuần lộc, Santa thực hiện
\texttt{prepareSleigh()} rồi phát 9 tín hiệu trên
\texttt{reindeerSem}.
\item Sau khi được giải phóng, từng tuần lộc thực hiện
\texttt{getHitched()}.
\end{itemize}

\paragraph{Quan hệ điều phối của chú lùn.}

\begin{itemize}
\item Semaphore \texttt{elfTex} đóng vai trò như cửa vào của nhóm
chú lùn hiện tại.
\item Chú lùn thứ nhất và thứ hai sau khi gia nhập nhóm sẽ mở lại
\texttt{elfTex}, cho phép chú lùn tiếp theo vào.
\item Chú lùn thứ ba đánh thức Santa nhưng không mở lại
\texttt{elfTex}; do đó các chú lùn đến sau bị chặn bên ngoài.
\item Sau khi nhóm đã đủ ba người, Santa thực hiện
\texttt{helpElves()} rồi phát ba tín hiệu trên \texttt{helpSem} để cho
ba chú lùn trong nhóm hiện tại thực hiện \texttt{getHelp()}.
\item Chỉ sau khi cả ba chú lùn hiện tại đã hoàn tất
\texttt{getHelp()}, chú lùn cuối cùng rời nhóm mới mở lại
\texttt{elfTex} cho nhóm kế tiếp.
\end{itemize}

Semaphore \texttt{helpSem} làm rõ quan hệ đồng bộ giữa Santa và nhóm
chú lùn: chú lùn không tự đi qua \texttt{getHelp()} ngay sau khi đánh
thức Santa, mà phải chờ Santa xử lý yêu cầu của nhóm hiện tại.

\subsubsection*{Giới hạn năng lực}

Bài toán có hai giới hạn bắt buộc:

\begin{itemize}
\item Santa chỉ chuẩn bị xe trượt tuyết cho một nhóm gồm 9 tuần lộc.
\item Chỉ đúng 3 chú lùn được thuộc nhóm xin trợ giúp hiện tại.
\end{itemize}

Khi chú lùn thứ ba gia nhập nhóm, \texttt{elfTex} không được mở lại cho
đến khi cả ba chú lùn đã thực hiện xong \texttt{getHelp()}.

\subsubsection*{Công bằng và ưu tiên}

Bài toán quy định tuần lộc có mức ưu tiên cao hơn chú lùn. Vì vậy, sau
khi được đánh thức, Santa kiểm tra yêu cầu của tuần lộc trước:

\begin{lstlisting}
if (reindeer >= 9) {
prepareSleigh();
}
else if (elves == 3) {
helpElves();
}
\end{lstlisting}

Cơ chế \texttt{elfTex} bảo đảm rằng một nhóm chú lùn đã hình thành sẽ
không bị nhóm mới chen vào trong khi đang nhận trợ giúp. Tuy nhiên, lời
giải vẫn giữ nguyên ưu tiên cho tuần lộc theo đặc tả của bài toán gốc.

% =========================================================

\subsection{Phân tích lỗi tiềm ẩn}

\subsubsection*{Race condition}

Giả sử hiện tại biến \texttt{elves = 1}. Hai chú lùn mới đến gần như
đồng thời và cùng muốn tham gia nhóm. Nếu không có mutex, có thể xảy ra
vết thực thi sau:

\begin{lstlisting}
Elf A doc elves = 1;
Elf B doc elves = 1;

Elf A ghi elves = 2;
Elf B ghi elves = 2;
\end{lstlisting}

Kết quả cuối cùng là \texttt{elves = 2}, mặc dù thực tế đã có đủ ba chú
lùn trong nhóm. Santa sẽ không được đánh thức đúng thời điểm.

Tương tự, nếu các tuần lộc cập nhật \texttt{reindeer} mà không có khóa,
hệ thống có thể ghi nhận sai số tuần lộc đã quay về.

\subsubsection*{Đoạn mã lỗi minh họa}

\begin{lstlisting}
// Giai phap sai: khong co mutex va khong khoa nhom elves
Elf_Naive() {
elves = elves + 1;

if (elves == 3) {
    signal(santaSem);
}

getHelp();

}
\end{lstlisting}

Đoạn mã trên thất bại vì biến \texttt{elves} không được bảo vệ bởi
mutex. Ngoài ra, không có \texttt{elfTex} để đóng cửa nhóm sau khi đủ
ba chú lùn; vì vậy chú lùn thứ tư có thể chen vào nhóm hiện tại.

\subsubsection*{Deadlock}

Một giải pháp sai có thể để chú lùn giữ \texttt{mutex} trong khi chờ
Santa thực hiện công việc:

\begin{lstlisting}
// Giai phap sai: chu lun giu mutex trong luc cho Santa
wait(mutex);                 // Khoa bien elves
elves = elves + 1;

if (elves == 3) {
signal(santaSem);        // Danh thuc Santa
}

wait(someSignal);            // Van giu mutex trong khi cho
signal(mutex);
\end{lstlisting}

Trong trường hợp này, Santa sau khi thức dậy cần lấy \texttt{mutex} để
kiểm tra biến \texttt{elves}. Tuy nhiên, chú lùn đang giữ chính mutex
đó và chờ Santa. Hai phía chờ nhau vô hạn, dẫn đến deadlock.

Lời giải đề xuất tránh tình huống trên vì chú lùn giải phóng
\texttt{mutex} trước khi thực hiện \texttt{getHelp()}, còn Santa chỉ
lấy mutex trong thời gian kiểm tra và xử lý điều kiện.

\subsubsection*{Starvation / fairness}

Tuần lộc được ưu tiên hơn chú lùn theo đúng yêu cầu bài toán. Do đó, nếu
hai loại yêu cầu cùng tồn tại, Santa sẽ xử lý tuần lộc trước.

Đối với các chú lùn, \texttt{elfTex} ngăn nhóm mới chen vào nhóm hiện
tại. Khi nhóm hiện tại đã hình thành, ba chú lùn trong nhóm đều được đi
qua \texttt{getHelp()} trước khi cửa được mở cho nhóm mới.

Lời giải không gây livelock, vì khi đủ điều kiện thì Santa hoặc nhóm
tiến trình tương ứng đều thực hiện được công việc thay vì liên tục
nhường quyền mà không tiến triển.

% =========================================================

\subsection{Giải pháp đề xuất}

\subsubsection*{Chiến lược tiếp cận}

Giải pháp trong phần này được trình bày lại dựa trên lời giải semaphore
của Downey trong mục \textit{Santa problem solution}
\cite[Sec.~5.5.2, pp.~141--142]{downey2016}. Giải pháp sử dụng:

\begin{itemize}
\item Một semaphore \texttt{mutex} để bảo vệ hai biến đếm
\texttt{elves} và \texttt{reindeer}.
\item Một semaphore \texttt{santaSem} để Santa ngủ và được đánh thức.
\item Một semaphore \texttt{reindeerSem} để tuần lộc chờ Santa chuẩn
bị xe.
\item Một semaphore \texttt{elfTex} để kiểm soát việc hình thành
từng nhóm 3 chú lùn.
\item Một semaphore \texttt{helpSem} để ba chú lùn trong nhóm hiện tại
chỉ thực hiện \texttt{getHelp()} sau khi Santa đã gọi
\texttt{helpElves()}.
\end{itemize}

Semaphore \texttt{elfTex} bảo đảm nhóm mới không thể gia nhập quá sớm,
còn \texttt{helpSem} bảo đảm nhóm hiện tại không nhận trợ giúp trước khi
Santa thật sự xử lý yêu cầu.

\subsection{Mã giả}

\subsubsection*{Khởi tạo cấu trúc dữ liệu và biến đồng bộ}

\begin{lstlisting}
int elves = 0;                // So chu lun trong nhom hien tai
int reindeer = 0;             // So tuan loc da quay ve va dang cho

semaphore santaSem = 0;       // Santa ngu va cho tin hieu danh thuc
semaphore reindeerSem = 0;    // Tuan loc cho Santa chuan bi xe
semaphore elfTex = 1;         // Cua vao cho tung nhom 3 chu lun
semaphore helpSem = 0;        // Santa cho phep tung chu lun nhan tro giup
semaphore mutex = 1;          // Bao ve hai bien elves va reindeer
\end{lstlisting}

\subsubsection*{Mã giả cho Santa}

\begin{lstlisting}
Santa() {
while (true) {
wait(santaSem);
// Cho den khi tuan loc thu 9 hoac chu lun thu 3 danh thuc Santa.

    wait(mutex);
    // Khoa hai bien dem de kiem tra dieu kien an toan.

    if (reindeer >= 9) {
        prepareSleigh();
        // Uu tien xu ly nhom tuan loc neu da co du 9 con.

        for (int i = 1; i <= 9; i++) {
            signal(reindeerSem);
            // Cho phep tung tuan loc thuc hien getHitched().
        }

        reindeer = reindeer - 9;
        // Loai nhom 9 tuan loc vua duoc xu ly khoi bien dem.
    }
    else if (elves == 3) {
        helpElves();
        // Santa giup nhom 3 chu lun hien tai.

        for (int i = 1; i <= 3; i++) {
            signal(helpSem);
            // Cho phep tung chu lun trong nhom goi getHelp().
        }
    }

    signal(mutex);
    // Mo khoa sau khi hoan tat viec kiem tra va xu ly.
}

}
\end{lstlisting}

\subsubsection*{Mã giả cho tuần lộc}

\begin{lstlisting}
Reindeer() {
while (true) {
returnFromVacation();
// Phan khong can dong bo: tuan loc quay ve sau ky nghi.

    wait(mutex);
    // Khoa bien reindeer truoc khi cap nhat.

    reindeer = reindeer + 1;
    // Ghi nhan mot tuan loc vua quay ve.

    if (reindeer == 9) {
        signal(santaSem);
        // Tuan loc thu 9 danh thuc Santa.
    }

    signal(mutex);
    // Mo khoa de cac tien trinh khac co the truy cap bien dem.

    wait(reindeerSem);
    // Cho Santa chuan bi xe truot tuyet.

    getHitched();
    // Tuan loc duoc mac vao xe.
}

}
\end{lstlisting}

\subsubsection*{Mã giả cho chú lùn}

\begin{lstlisting}
Elf() {
while (true) {
makeToys();
// Phan khong can dong bo: chu lun lam do choi.

    encounterProblem();
    // Chu lun gap kho khan va can tro giup.

    wait(elfTex);
    // Xin quyen gia nhap nhom hien tai.
    // Neu nhom da du 3 nguoi, chu lun moi phai cho tai day.

    wait(mutex);
    // Khoa bien elves truoc khi cap nhat.

    elves = elves + 1;
    // Them mot chu lun vao nhom xin tro giup.

    if (elves == 3) {
        signal(santaSem);
        // Chu lun thu 3 danh thuc Santa.
        // Khong mo elfTex vi nhom hien tai da du ba nguoi.
    }
    else {
        signal(elfTex);
        // Nhom chua du ba nguoi, mo cua cho chu lun ke tiep.
    }

    signal(mutex);
    // Mo khoa sau khi cap nhat bien elves.

    wait(helpSem);
    // Cho den khi Santa thuc su xu ly nhom hien tai.

    getHelp();
    // Chu lun thuc hien viec nhan tro giup.

    wait(mutex);
    // Khoa bien elves truoc khi roi nhom hien tai.

    elves = elves - 1;
    // Ghi nhan chu lun da hoan tat getHelp().

    if (elves == 0) {
        signal(elfTex);
        // Chu lun cuoi cung mo cua cho nhom tiep theo.
    }

    signal(mutex);
    // Mo khoa sau khi cap nhat xong.
}

}
\end{lstlisting}

\subsection{Kiểm chứng ngắn}

Giả sử lần lượt có ba chú lùn gặp khó khăn. Chú lùn thứ nhất và thứ hai
tăng biến \texttt{elves}, sau đó mở lại \texttt{elfTex} để cho phép chú
lùn kế tiếp gia nhập nhóm. Khi chú lùn thứ ba đi vào,
\texttt{elves = 3}; chú lùn này đánh thức Santa nhưng không mở lại
\texttt{elfTex}. Vì vậy, chú lùn thứ tư không thể chen vào nhóm hiện tại.

Sau khi nhóm đã đủ ba người, các chú lùn chờ trên \texttt{helpSem}.
Santa thức dậy, gọi \texttt{helpElves()} và phát ba tín hiệu
\texttt{helpSem}; khi đó ba chú lùn trong nhóm hiện tại mới thực hiện
\texttt{getHelp()}. Mỗi chú lùn sau khi hoàn tất \texttt{getHelp()} sẽ
giảm biến \texttt{elves} trong vùng được bảo vệ bởi \texttt{mutex}. Chú lùn cuối cùng làm
\texttt{elves = 0} rồi mở lại \texttt{elfTex}; từ thời điểm đó nhóm mới
mới có thể được hình thành.

Nếu tuần lộc thứ 9 quay về, tiến trình này đánh thức Santa. Khi Santa
kiểm tra thấy \texttt{reindeer >= 9}, Santa ưu tiên gọi
\texttt{prepareSleigh()} và giải phóng đúng 9 tuần lộc thông qua
\texttt{reindeerSem}. Mọi cập nhật biến đếm đều nằm trong vùng bảo vệ
của \texttt{mutex}, nhờ đó tránh được race condition.

\subsection{Kết luận bài toán}

Bài toán Ông già Noel minh họa việc đồng bộ nhiều nhóm tiến trình với
điều kiện kích hoạt và mức ưu tiên khác nhau. Thông qua
\texttt{mutex}, \texttt{santaSem}, \texttt{reindeerSem} và
\texttt{elfTex}, lời giải bảo đảm Santa xử lý đúng điều kiện, tuần lộc
được ưu tiên theo đặc tả và các chú lùn được chia thành từng nhóm đúng
ba người.

\section{Bài toán Trông trẻ (Child Care Problem)}

Bài toán Child Care đặt ra ràng buộc tỷ lệ giữa số trẻ em và số người lớn trong phòng. Trẻ em chỉ được vào nếu sau khi vào vẫn thỏa mãn tỷ lệ an toàn; người lớn chỉ được rời đi nếu sau khi rời đi tỷ lệ vẫn hợp lệ.

\subsection{Mô tả bài toán}
Trong một khu trông trẻ có hai loại tiến trình: \textit{adult} (người lớn) và \textit{child} (trẻ em). Ràng buộc an toàn yêu cầu số trẻ em trong phòng luôn không vượt quá một bội số của số người lớn. Biến thể phổ biến đặt tỷ lệ:
\[
    children \le 3 \cdot adults.
\]
Mục tiêu: đồng bộ để mọi thao tác vào/ra đều giữ đúng tỷ lệ, đồng thời hệ thống vẫn có progress (không tự tạo deadlock do thiết kế sai điều kiện chờ).

\subsection{Các thực thể tham gia}
\begin{itemize}
    \item \textbf{Adult:} tiến trình người lớn, có thể vào phòng hoặc xin rời phòng.
    \item \textbf{Child:} tiến trình trẻ em, có thể vào phòng nếu tỷ lệ an toàn vẫn được giữ và có thể rời phòng bất kỳ lúc nào.
    \item \textbf{Monitor phòng trông trẻ:} thành phần bảo vệ biến đếm và điều kiện chờ để mọi thao tác vào/ra được kiểm tra nguyên tử.
\end{itemize}

\subsection{Tài nguyên dùng chung}
\begin{itemize}
    \item Biến đếm $adults$ và $children$.
    \item Hàng đợi chờ của trẻ muốn vào và người lớn muốn rời.
\end{itemize}

\subsection{Ràng buộc đồng bộ}
\begin{itemize}
    \item \textbf{Trẻ vào:} chỉ được vào nếu sau khi vào vẫn thỏa $children + 1 \le 3 \cdot adults$.
    \item \textbf{Người lớn rời:} chỉ được rời nếu sau khi rời vẫn thỏa $children \le 3 \cdot (adults - 1)$.
    \item \textbf{Trẻ rời / người lớn vào} luôn hợp lệ (vì làm giảm hoặc nới ràng buộc), nhưng vẫn cần khóa để cập nhật biến đếm an toàn.
\end{itemize}

\subsection{Phân tích lỗi tiềm ẩn}

\subsubsection*{Race condition}
Nếu cập nhật \texttt{adults} và \texttt{children} ngoài monitor, hai tiến trình
có thể cùng kiểm tra điều kiện trên dữ liệu cũ và làm sai bất biến
$children \le 3 \cdot adults$.

\subsubsection*{Deadlock}
Người lớn rời phòng phải chờ bằng condition variable của monitor, tức là thao
tác \texttt{wait} tạm nhả mutex. Nếu chờ mà vẫn giữ mutex, người lớn khác hoặc
trẻ rời phòng không thể cập nhật trạng thái để làm điều kiện trở thành đúng.

\subsubsection*{Starvation / fairness}
Nếu quên đánh thức sau khi người lớn vào hoặc trẻ rời, tiến trình đang chờ có
thể ngủ vô hạn dù điều kiện đã hợp lệ. Ngoài ra, luôn dùng \texttt{while} khi
chờ để chống đánh thức giả và lịch trình xen kẽ bất lợi.

\subsection{Giải pháp đề xuất}
Cách trình bày rõ ràng là dùng monitor (mutex + condition variable):
\begin{itemize}
    \item Khóa $mutex$ bảo vệ cập nhật $adults, children$.
    \item $canChildEnter$ để trẻ chờ khi tỷ lệ chưa đủ an toàn.
    \item $canAdultLeave$ để người lớn chờ khi rời đi sẽ làm tỷ lệ bị vi phạm.
\end{itemize}
Sau khi một tiến trình vào/ra, cần \textit{signal/broadcast} để đánh thức các tiến trình đang chờ vì điều kiện có thể đã thay đổi.

\subsection{Mã giả}

\captionof{algorithm}{Child Care (tỷ lệ $children \le 3 \cdot adults$) -- giả mã monitor}
\label{alg:child-care}
\begin{algorithmic}[1]
\State \textbf{Shared:} \texttt{mutex}; condition \texttt{canChildEnter}, \texttt{canAdultLeave}
\State \textbf{Shared:} \texttt{adults} $\leftarrow 0$; \texttt{children} $\leftarrow 0$

\Procedure{AdultEnter}{}
    \State \Call{lock}{\texttt{mutex}}
    \State $\texttt{adults} \gets \texttt{adults} + 1$
    \State \Call{broadcast}{\texttt{canChildEnter}}
    \State \Call{broadcast}{\texttt{canAdultLeave}}
    \State \Call{unlock}{\texttt{mutex}}
\EndProcedure

\Procedure{AdultLeave}{}
    \State \Call{lock}{\texttt{mutex}}
    \While{$\texttt{children} > 3 \cdot (\texttt{adults} - 1)$}
        \State \Call{wait}{\texttt{canAdultLeave}, \texttt{mutex}}
    \EndWhile
    \State $\texttt{adults} \gets \texttt{adults} - 1$
    \State \Call{broadcast}{\texttt{canChildEnter}}
    \State \Call{broadcast}{\texttt{canAdultLeave}}
    \State \Call{unlock}{\texttt{mutex}}
\EndProcedure

\Procedure{ChildEnter}{}
    \State \Call{lock}{\texttt{mutex}}
    \While{$\texttt{children} + 1 > 3 \cdot \texttt{adults}$}
        \State \Call{wait}{\texttt{canChildEnter}, \texttt{mutex}}
    \EndWhile
    \State $\texttt{children} \gets \texttt{children} + 1$
    \State \Call{unlock}{\texttt{mutex}}
\EndProcedure

\Procedure{ChildLeave}{}
    \State \Call{lock}{\texttt{mutex}}
    \State $\texttt{children} \gets \texttt{children} - 1$
    \State \Call{broadcast}{\texttt{canChildEnter}}
    \State \Call{broadcast}{\texttt{canAdultLeave}}
    \State \Call{unlock}{\texttt{mutex}}
\EndProcedure
\end{algorithmic}

\subsection{Kiểm chứng ngắn}

Mọi thay đổi lên \texttt{adults} và \texttt{children} đều diễn ra khi giữ \texttt{mutex}, nên không có hai tiến trình đồng thời cập nhật sai biến đếm. Trẻ chỉ vào khi điều kiện sau khi vào vẫn thỏa $children + 1 \le 3 \cdot adults$; người lớn chỉ rời khi điều kiện sau khi rời vẫn thỏa $children \le 3 \cdot (adults - 1)$. Khi người lớn vào, cả trẻ đang chờ và người lớn đang chờ rời đi đều có thể được đánh thức vì mẫu số an toàn tăng lên. Khi trẻ rời, cả hai condition variable cũng được đánh thức vì việc giảm \texttt{children} có thể đồng thời cho phép trẻ khác vào hoặc người lớn rời đi.

\subsection{Kết luận bài toán}

Child Care là bài toán monitor điển hình với bất biến an toàn $children \le 3 \cdot adults$. Lời giải đúng cần kiểm tra điều kiện bằng vòng lặp, cập nhật biến đếm trong vùng khóa và đánh thức lại các hàng đợi liên quan sau mỗi thay đổi có thể làm điều kiện chờ trở thành đúng.

\section{Bài toán Ghép đội đấu hạng (Multiplayer Matchmaking Problem)}

\subsection{Mô tả bài toán}

\subsubsection*{Cốt truyện và bối cảnh}

Bài toán Ghép đội đấu hạng mô phỏng một hệ thống ghép trận trong trò chơi
nhiều người chơi trực tuyến. Mỗi người chơi gửi yêu cầu tìm trận, mang theo
một số thuộc tính như điểm xếp hạng, vai trò mong muốn, khu vực máy chủ và thời
điểm bắt đầu chờ. Hệ thống phải chọn một nhóm người chơi đủ điều kiện để tạo
thành một trận đấu hợp lệ.

Trong báo cáo này, ta xét một biến thể 5v5 có ba vai trò cơ bản:
\textit{Tank}, \textit{Damage} và \textit{Support}. Mỗi đội cần đúng:

\[
    1 \text{ Tank} + 3 \text{ Damage} + 1 \text{ Support}.
\]

Vì một trận gồm hai đội, hệ thống cần chọn tổng cộng:

\[
    2 \text{ Tank} + 6 \text{ Damage} + 2 \text{ Support}.
\]

Ngoài ràng buộc số lượng và vai trò, người chơi trong cùng trận chỉ được ghép
nếu điểm xếp hạng của họ nằm trong một khoảng chênh lệch cho phép. Khoảng này
có thể mở rộng dần theo thời gian chờ để tránh người chơi bị kẹt quá lâu.

Bài toán này không phải là một bài toán semaphore kinh điển với một lời giải
duy nhất. Nó là một mô hình thực tế dùng để áp dụng các nguyên lý đồng bộ:
bảo vệ hàng đợi dùng chung, chọn nhóm một cách nguyên tử, đánh thức đúng người
được chọn và bảo đảm những người không được chọn tiếp tục chờ trong trạng thái
nhất quán.

\subsection{Các thực thể tham gia}

\paragraph{Người chơi (Player Thread).}
Mỗi người chơi là một tiến trình hoặc luồng gửi yêu cầu tìm trận. Người chơi
cung cấp thông tin về rank, vai trò và khu vực, sau đó chờ đến khi được hệ
thống đưa vào một trận cụ thể.

\begin{center}
\texttt{Gửi yêu cầu $\rightarrow$ Vào hàng đợi $\rightarrow$ Chờ ghép trận
$\rightarrow$ Nhận matchId $\rightarrow$ Vào phòng đấu}
\end{center}

\paragraph{Bộ ghép trận (Matchmaker).}
Bộ ghép trận là tiến trình điều phối. Nó quan sát các hàng đợi người chơi, tìm
một nhóm thỏa điều kiện, lấy nhóm đó ra khỏi hàng đợi và đánh thức đúng các
người chơi được chọn.

\begin{center}
\texttt{Chờ yêu cầu $\rightarrow$ Quét hàng đợi $\rightarrow$ Chọn nhóm hợp lệ
$\rightarrow$ Tạo trận $\rightarrow$ Đánh thức người chơi}
\end{center}

\paragraph{Phòng đấu (Match Room).}
Phòng đấu là kết quả sau khi ghép thành công. Mỗi phòng có một định danh
\texttt{matchId}, danh sách người chơi và thông tin đội. Sau khi một người chơi
được gán vào phòng, yêu cầu của người đó không còn nằm trong hàng đợi tìm trận.

\subsection{Tài nguyên dùng chung}

Các tài nguyên dùng chung trong bài toán gồm:

\begin{itemize}
    \item Các hàng đợi \texttt{queue[role]} chứa người chơi đang chờ theo vai trò.
    \item Trạng thái từng yêu cầu: đang chờ, đã được chọn, đã hủy hoặc đã vào trận.
    \item Bộ đếm và danh sách người chơi đang chờ theo rank/role/khu vực.
    \item Bộ sinh \texttt{matchId} cho các trận mới.
    \item Semaphore hoặc condition variable riêng của từng người chơi để chờ
    kết quả ghép trận.
\end{itemize}

Tất cả cấu trúc này phải được bảo vệ bởi \texttt{mutex} hoặc nằm bên trong một
monitor. Nếu không, hai tiến trình matchmaker có thể cùng chọn một người chơi
cho hai trận khác nhau, hoặc một người chơi vừa hủy tìm trận vẫn bị đưa vào
phòng đấu.

\subsection{Ràng buộc đồng bộ}

\subsubsection*{Loại trừ tương hỗ}

Đoạn găng của bài toán nằm ở các thao tác:

\begin{itemize}
    \item Thêm yêu cầu mới vào hàng đợi.
    \item Xóa yêu cầu khi người chơi hủy tìm trận.
    \item Quét và chọn nhóm người chơi hợp lệ.
    \item Gán \texttt{matchId} cho nhóm đã chọn.
    \item Đánh dấu trạng thái người chơi từ \texttt{WAITING} sang
    \texttt{MATCHED}.
\end{itemize}

Đặc biệt, thao tác chọn nhóm và xóa nhóm khỏi hàng đợi phải là một hành động
nguyên tử. Nếu không, hai matchmaker chạy đồng thời có thể cùng nhìn thấy một
nhóm hợp lệ và cùng sử dụng một vài người chơi trùng nhau.

\subsubsection*{Đồng bộ điều kiện}

Người chơi sau khi vào hàng đợi phải chờ trên semaphore riêng
\texttt{matchedSem}. Người chơi chỉ được tiếp tục khi một matchmaker đã:

\begin{itemize}
    \item Chọn người chơi đó vào một nhóm hợp lệ.
    \item Gán \texttt{matchId}.
    \item Xác định đội và vai trò trong trận.
    \item Xóa yêu cầu khỏi hàng đợi chờ.
\end{itemize}

Matchmaker cũng có thể ngủ trên semaphore \texttt{requestSem} khi toàn bộ hệ
thống chưa có đủ yêu cầu. Mỗi người chơi mới đến sẽ phát tín hiệu để đánh thức
matchmaker kiểm tra lại khả năng tạo trận. Vì ngưỡng rank được mở rộng theo
thời gian chờ, matchmaker không chỉ phụ thuộc vào yêu cầu mới mà còn cần thức
dậy định kỳ để quét lại hàng đợi.

\subsubsection*{Giới hạn năng lực}

Trong biến thể 5v5, mỗi trận phải có đúng 10 người chơi:

\[
    \texttt{MATCH\_SIZE} = 10.
\]

Với cấu hình vai trò đã chọn, điều kiện về số lượng là:

\[
    Tank \geq 2,\quad Damage \geq 6,\quad Support \geq 2.
\]

Ngoài ra, mỗi người chơi chỉ được thuộc về tối đa một trận tại một thời điểm.
Đây là bất biến quan trọng nhất:

\[
    \forall p,\quad p \text{ không thể xuất hiện trong hai } matchId
    \text{ khác nhau.}
\]

\subsubsection*{Ràng buộc chất lượng trận đấu}

Một trận hợp lệ không chỉ đủ người. Hệ thống còn cần kiểm tra độ lệch rank:

\[
    \max(MMR) - \min(MMR) \leq \Delta(t).
\]

Ở đây \(\Delta(t)\) là ngưỡng chênh lệch cho phép, có thể tăng theo thời gian
chờ \(t\). Khi người chơi mới vào hàng đợi, \(\Delta(t)\) nhỏ để ưu tiên chất
lượng trận. Nếu người chơi chờ lâu, ngưỡng được nới dần để tăng khả năng tìm
trận.

\subsubsection*{Công bằng và ưu tiên}

Nếu hệ thống luôn chọn nhóm có rank gần nhau nhất, một số người chơi ở rank
quá cao, quá thấp hoặc vai trò hiếm có thể chờ rất lâu. Vì vậy, lời giải cần
có cơ chế \textit{aging}: yêu cầu càng chờ lâu thì càng được ưu tiên, hoặc
ngưỡng rank càng được mở rộng.

Một chính sách hợp lý là:

\begin{itemize}
    \item Ưu tiên người chơi chờ lâu trong từng hàng đợi vai trò.
    \item Không lấy người chơi mới đến để thay thế người đã chờ lâu nếu cả hai
    đều tạo được trận hợp lệ.
    \item Mở rộng khoảng rank theo thời gian chờ nhưng vẫn giữ một ngưỡng tối đa
    để tránh trận quá lệch.
\end{itemize}

\subsection{Phân tích lỗi tiềm ẩn}

\subsubsection*{Race condition}

Giả sử có hai matchmaker chạy song song để tăng tốc. Nếu cả hai cùng đọc hàng
đợi mà không khóa, có thể xảy ra vết thực thi sau:

\begin{lstlisting}
Matchmaker A doc thay du 2 Tank, 6 Damage, 2 Support;
Matchmaker B cung doc thay dung nhom nguoi cho do;

Matchmaker A tao matchId = 101 cho 10 nguoi nay;
Matchmaker B tao matchId = 102 cho cung mot vai nguoi trong nhom;
\end{lstlisting}

Kết quả là một người chơi có thể nhận hai thông báo vào trận khác nhau. Đây là
lỗi nghiêm trọng vì trạng thái hàng đợi và trạng thái phòng đấu không còn nhất
quán.

\subsubsection*{Đánh thức sai người chơi}

Nếu matchmaker dùng một semaphore chung cho tất cả người chơi, một người không
thuộc nhóm được chọn vẫn có thể bị đánh thức. Người này có thể đọc
\texttt{matchId} rỗng hoặc đi vào nhánh lỗi.

\begin{lstlisting}
// Giai phap sai: semaphore chung khong gan voi tung player
signal(globalMatched);
\end{lstlisting}

Lời giải đúng nên dùng semaphore riêng trong từng \texttt{PlayerRequest}. Khi
đó matchmaker phát tín hiệu đúng vào các yêu cầu đã được gán \texttt{matchId}.

\subsubsection*{Hủy tìm trận trong lúc đang ghép}

Người chơi có thể hủy tìm trận. Nếu thao tác hủy không đồng bộ với thao tác
ghép, matchmaker có thể chọn một yêu cầu vừa bị hủy. Do đó trạng thái yêu cầu
cần được kiểm tra trong cùng đoạn găng:

\begin{itemize}
    \item Nếu request còn \texttt{WAITING}, matchmaker được quyền chọn.
    \item Nếu request đã \texttt{CANCELLED}, matchmaker phải bỏ qua.
    \item Khi matchmaker chuyển request sang \texttt{MATCHED}, thao tác hủy sau
    đó không được xóa request khỏi trận đã tạo.
\end{itemize}

\subsubsection*{Deadlock}

Deadlock có thể xảy ra nếu matchmaker giữ \texttt{mutex} rồi chờ người chơi xác
nhận vào phòng, trong khi người chơi cần lấy cùng \texttt{mutex} để đọc trạng
thái trận:

\begin{lstlisting}
// Giai phap sai: giu mutex trong khi cho player phan hoi
wait(mutex);
signal(player.matchedSem);
wait(player.acceptedSem);
signal(mutex);
\end{lstlisting}

Nếu người chơi sau khi thức dậy cần \texttt{mutex} để đọc \texttt{matchId} và
phát \texttt{acceptedSem}, hai bên sẽ chờ nhau. Cách tránh là matchmaker chỉ
giữ mutex trong lúc cập nhật trạng thái, sau đó nhả mutex rồi mới chờ các xác
nhận ngoài vùng găng nếu hệ thống cần pha xác nhận.

\subsubsection*{Starvation / fairness}

Starvation xuất hiện khi một nhóm người chơi liên tục bị bỏ qua vì không tạo
được trận chất lượng cao theo tiêu chí ban đầu. Ví dụ, vai trò Tank đang thiếu
có thể khiến nhiều Damage chờ lâu; hoặc một người chơi rank rất cao khó tìm đủ
đối thủ cùng mức. Cơ chế aging và mở rộng \(\Delta(t)\) giúp giảm rủi ro này.

Livelock có thể xảy ra nếu hệ thống liên tục thử tạo trận, thấy không đạt tiêu
chí, thay đổi ứng viên rồi lại quay về trạng thái cũ mà không ai được ghép.
Để tránh livelock, mỗi lần quét nên có quy tắc chọn ứng viên ổn định, ví dụ lấy
người chờ lâu nhất làm hạt nhân rồi mở rộng dần tập ứng viên quanh người đó.

\subsection{Giải pháp đề xuất}

\subsubsection*{Chiến lược tiếp cận}

Lời giải sử dụng monitor cho hàng đợi matchmaking. Monitor gom các hàng đợi,
bộ sinh \texttt{matchId}, các thủ tục thêm/hủy yêu cầu và thủ tục tạo trận.
Mỗi người chơi có một semaphore riêng để chờ kết quả.

Ý tưởng chính:

\begin{itemize}
    \item Người chơi tạo \texttt{PlayerRequest}, đưa vào hàng đợi theo vai trò
    rồi ngủ trên semaphore riêng.
    \item Matchmaker lấy mutex, tìm một nhóm hợp lệ theo vai trò và rank.
    \item Nếu tìm được nhóm, matchmaker xóa nhóm khỏi hàng đợi, gán
    \texttt{matchId}, đổi trạng thái thành \texttt{MATCHED}, rồi đánh thức đúng
    10 người chơi.
    \item Nếu chưa tìm được nhóm, matchmaker chờ yêu cầu mới hoặc tự thức dậy
    theo chu kỳ ngắn để áp dụng chính sách aging.
\end{itemize}

\subsection{Mã giả}

\subsubsection*{Khởi tạo cấu trúc dữ liệu và biến đồng bộ}

\begin{lstlisting}
const int TEAM_SIZE = 5;
const int MATCH_SIZE = 10;

enum Role { TANK, DAMAGE, SUPPORT };
enum Status { WAITING, MATCHED, CANCELLED };

typedef struct PlayerRequest {
    int playerId;
    int mmr;
    Role role;
    long enqueueTime;

    Status status;
    int matchId;
    int teamId;

    semaphore matchedSem;
} PlayerRequest;

Queue queue[3];              // Hang doi theo role
int nextMatchId = 1;

semaphore mutex = 1;         // Bao ve queue va trang thai request
semaphore requestSem = 0;    // Bao co request moi
const int RESCAN_INTERVAL_MS = 1000; // Chu ky quet lai de ap dung aging
\end{lstlisting}

\subsubsection*{Hàm kiểm tra nhóm hợp lệ}

\begin{lstlisting}
bool isValidGroup(List group, long now) {
    if (countRole(group, TANK) != 2) return false;
    if (countRole(group, DAMAGE) != 6) return false;
    if (countRole(group, SUPPORT) != 2) return false;

    int minMmr = minMMR(group);
    int maxMmr = maxMMR(group);
    long oldestWait = now - minEnqueueTime(group);

    int delta = allowedDelta(oldestWait);

    return (maxMmr - minMmr <= delta);
}
\end{lstlisting}

Hàm \texttt{allowedDelta} biểu diễn chính sách mở rộng khoảng rank theo thời
gian chờ. Ví dụ, hệ thống có thể bắt đầu với khoảng hẹp rồi tăng dần đến một
ngưỡng tối đa.

\subsubsection*{Mã giả cho người chơi}

\begin{lstlisting}
PlayerThread(int playerId, int mmr, Role role) {
    PlayerRequest self;

    self.playerId = playerId;
    self.mmr = mmr;
    self.role = role;
    self.enqueueTime = now();
    self.status = WAITING;
    self.matchId = -1;
    self.teamId = -1;
    self.matchedSem = 0;

    wait(mutex);
    queue[role].enqueue(&self);
    signal(mutex);

    signal(requestSem);
    // Bao cho matchmaker rang co yeu cau moi.

    wait(self.matchedSem);
    // Cho den khi duoc gan vao mot tran cu the.

    if (self.status == MATCHED) {
        joinMatch(self.matchId, self.teamId);
    }
}
\end{lstlisting}

\subsubsection*{Mã giả cho hủy tìm trận}

\begin{lstlisting}
CancelSearch(PlayerRequest *self) {
    wait(mutex);

    if (self->status == WAITING) {
        self->status = CANCELLED;
        removeFromQueue(queue[self->role], self);
    }

    signal(mutex);
}
\end{lstlisting}

Thủ tục hủy chỉ có hiệu lực khi yêu cầu còn đang chờ. Nếu matchmaker đã chuyển
trạng thái sang \texttt{MATCHED}, yêu cầu không còn thuộc hàng đợi và người chơi
đã được gán vào một trận cụ thể.

\subsubsection*{Mã giả cho matchmaker}

\begin{lstlisting}
MatchmakerThread() {
    while (true) {
        timedWait(requestSem, RESCAN_INTERVAL_MS);
        // Thuc day khi co request moi hoac khi timer het han.

        wait(mutex);

        while (true) {
            List group = findCandidateGroup(queue, now());

            if (group.empty()) {
                break;
            }

            if (!isValidGroup(group, now())) {
                break;
            }

            int id = nextMatchId;
            nextMatchId = nextMatchId + 1;

            removeGroupFromQueues(group);
            assignTeams(group);

            for (PlayerRequest *p in group) {
                p->status = MATCHED;
                p->matchId = id;
                signal(p->matchedSem);
            }
        }

        signal(mutex);
    }
}
\end{lstlisting}

Trong đoạn mã trên, thao tác tìm nhóm, xóa nhóm khỏi hàng đợi và gán
\texttt{matchId} diễn ra trong cùng đoạn găng. Vì vậy, không có matchmaker thứ
hai nào có thể lấy trùng người chơi trong lúc nhóm đang được tạo.

\subsubsection*{Hàm chọn nhóm ứng viên}

\begin{lstlisting}
List findCandidateGroup(Queue queue[], long now) {
    PlayerRequest *seed = oldestWaitingPlayer(queue);

    if (seed == null) {
        return emptyList();
    }

    int delta = allowedDelta(now - seed->enqueueTime);

    List group;
    group.add(seed);
    // Bat buoc dua nguoi cho lau nhat vao nhom ung vien.

    int needTank = 2;
    int needDamage = 6;
    int needSupport = 2;

    if (seed->role == TANK) needTank = needTank - 1;
    if (seed->role == DAMAGE) needDamage = needDamage - 1;
    if (seed->role == SUPPORT) needSupport = needSupport - 1;

    group.add(selectOldestCompatibleExcept(queue[TANK], needTank,
                                           seed, seed->mmr, delta));
    group.add(selectOldestCompatibleExcept(queue[DAMAGE], needDamage,
                                           seed, seed->mmr, delta));
    group.add(selectOldestCompatibleExcept(queue[SUPPORT], needSupport,
                                           seed, seed->mmr, delta));

    if (group.size() != MATCH_SIZE) {
        return emptyList();
    }

    return group;
}
\end{lstlisting}

Cách chọn người chờ lâu nhất làm \textit{seed} giúp hệ thống không chỉ tối ưu
rank gần nhau mà còn quan tâm đến thời gian chờ. Seed được đưa vào nhóm trước,
sau đó hệ thống mới lấy thêm các người chơi tương thích theo từng vai trò. Nếu
không đủ người tương thích quanh seed, hàm trả về rỗng và matchmaker sẽ thử lại
ở lần quét sau, khi có yêu cầu mới hoặc khi ngưỡng aging được mở rộng.

\subsection{Kiểm chứng ngắn}

Giả sử hàng đợi hiện có 2 Tank, 6 Damage và 2 Support tương thích về rank.
Matchmaker lấy \texttt{mutex}, gọi \texttt{findCandidateGroup} và nhận đúng 10
yêu cầu. Sau đó matchmaker xóa cả nhóm khỏi các hàng đợi, tạo
\texttt{matchId}, chia hai đội và gán trạng thái \texttt{MATCHED}. Chỉ sau khi
trạng thái của từng người đã nhất quán, matchmaker mới phát tín hiệu trên
\texttt{matchedSem} riêng của họ.

Nếu một người chơi khác đến trong lúc matchmaker đang tạo trận, người đó bị
chặn tại \texttt{wait(mutex)} và không thể chen vào nhóm đang được chốt. Nếu
một người chơi cố hủy tìm trận sau khi đã được gán \texttt{MATCHED}, thủ tục
hủy sẽ không xóa người đó khỏi trận vì trạng thái không còn là
\texttt{WAITING}. Nhờ vậy, hệ thống giữ được bất biến mỗi người chơi thuộc tối
đa một trận.

\subsection{Kết luận bài toán}

Bài toán Ghép đội đấu hạng cho thấy đồng bộ tài nguyên găng không chỉ xuất hiện
trong các ví dụ kinh điển mà còn trong hệ thống dịch vụ thời gian thực. Tài
nguyên găng ở đây là hàng đợi người chơi và trạng thái yêu cầu tìm trận. Lời
giải đúng phải chọn nhóm một cách nguyên tử, đánh thức đúng người chơi được
chọn, xử lý hủy tìm trận an toàn và có chính sách fairness để tránh starvation.
Monitor kết hợp semaphore riêng cho từng người chơi là một mô hình phù hợp vì
nó tách rõ hai việc: bảo vệ trạng thái dùng chung và thông báo kết quả ghép
trận cho từng tiến trình.
